\section{The AKLT state}\label{sec:aklt}
%% ===================================================================

The Affleck--Kennedy--Lieb--Tasaki (AKLT)
state~\cite{AKLT1987Rigorous,AKLT1988ValenceBond} is the canonical
example of a normal (block-injective) MPS with non-trivial
symmetry-protected topological (SPT)
order~\cite{PerezGarcia2008StringOrder,Schuch2011Phases,Cirac2021Matrix}.

\begin{definition}[AKLT tensor from singlet bonds]\label{def:aklt_tensor_review}
    \lean{MPSTensor.rmpSingletY, MPSTensor.akltRMPSpinOneState, MPSTensor.akltTensorRMP}
    \leanok
    \uses{def:mps_tensor}
    Place a spin-$\tfrac12$ singlet on every bond, encoded by
    \begin{align}
        Y & = \begin{pmatrix}0&-1\\1&0\end{pmatrix},
        \notag
    \end{align}
    and project the two spin-$\tfrac12$ particles at each site onto their joint
    spin-$1$ subspace. With the physical states labelled by $S_z=0,\pm1$, the
    resulting tensor with $d=3$ and $D=2$ is
    \begin{align}
        A^{+1} & = \begin{pmatrix}1&0\\0&0\end{pmatrix}Y,
               & A^{0}                                    & = \frac{1}{\sqrt2}\begin{pmatrix}0&1\\1&0\end{pmatrix}Y,
               & A^{-1}                                   & = \begin{pmatrix}0&0\\0&1\end{pmatrix}Y.
        \label{eq:aklt_review_tensor}
    \end{align}
    The three left factors are the spin-$1$ states $\ket{\uparrow\uparrow}$,
    $(\ket{\uparrow\downarrow}+\ket{\downarrow\uparrow})/\sqrt2$, and
    $\ket{\downarrow\downarrow}$ written as coefficient matrices. This is the
    tensor of \cite[Appendix~A, ``The AKLT state'']{Cirac2021Matrix}, for the
    state of~\cite{AKLT1987Rigorous}.
\end{definition}

\begin{lemma}[Matrices of the singlet-bond AKLT tensor]\label{lem:aklt_review_matrices}
    \lean{MPSTensor.akltTensorRMP_zero, MPSTensor.akltTensorRMP_one,
        MPSTensor.akltTensorRMP_two}
    \leanok
    \uses{def:aklt_tensor_review}
    The matrices of (\ref{eq:aklt_review_tensor}) are
    $A^{0}=\sigma_z/\sqrt2$, $A^{+1}=-\ketbra{0}{1}$, and
    $A^{-1}=\ketbra{1}{0}$.
\end{lemma}

\begin{theorem}[The site map projects onto spin $1$]\label{thm:aklt_review_spin_one}
    \lean{MPSTensor.akltRMPSpinOneState_transpose,
        MPSTensor.akltRMPSpinOneState_orthonormal}
    \leanok
    \uses{def:aklt_tensor_review}
    The three coefficient matrices $P^{s}$ with $A^{s}=P^{s}Y$ are symmetric,
    $(P^{s})^{T}=P^{s}$, so they are states of the symmetric (spin-$1$)
    subspace of two spin-$\tfrac12$ particles, and they are orthonormal:
    $\tr\big((P^{s})^{\dagger}P^{t}\big)=\delta_{st}$.
\end{theorem}

\begin{definition}[The Pauli basis of the spin-$1$ site]\label{def:aklt_review_basis}
    \lean{MPSTensor.akltRMPBasis, MPSTensor.akltRMPBasisFixed}
    \leanok
    \uses{def:aklt_tensor_review}
    The basis of \cite[Appendix~A, ``The AKLT state'']{Cirac2021Matrix} is
    \begin{align}
        \ket{+} & = \frac{i}{\sqrt2}\big(\ket{-1}+\ket{+1}\big),
                & \ket{-}                                        & = \frac{1}{\sqrt2}\big(\ket{-1}-\ket{+1}\big),
                & \ket{0}.
        \label{eq:aklt_review_basis}
    \end{align}
    The second basis replaces $\ket{+}$ with
    $-\frac{i}{\sqrt2}(\ket{-1}+\ket{+1})$ and keeps $\ket{-}$ and $\ket{0}$.
\end{definition}

\begin{theorem}[The Pauli bases are orthonormal]\label{thm:aklt_review_basis_unitary}
    \lean{MPSTensor.akltRMPBasis_mem_unitaryGroup,
        MPSTensor.akltRMPBasisFixed_mem_unitaryGroup}
    \leanok
    \uses{def:aklt_review_basis}
    The basis (\ref{eq:aklt_review_basis}) is orthonormal, and so is the basis
    obtained by replacing $\ket{+}$ with $-\frac{i}{\sqrt2}(\ket{-1}+\ket{+1})$:
    the matrices whose columns are these vectors are unitary.
\end{theorem}
\begin{proof}
    \leanok
    The vectors $\ket{\pm}$ are supported on $\ket{\pm1}$ and $\ket{0}$ is
    orthogonal to both; the inner product of $\ket{+}$ and $\ket{-}$ is
    $\frac{-i}{2}(1-1)=0$, and each vector has norm $1$. A phase on $\ket{+}$
    preserves all of these relations.
\end{proof}

\begin{theorem}[Pauli form of the AKLT tensor]\label{thm:aklt_review_pauli_form}
    \lean{MPSTensor.akltTensorRMPPauli, MPSTensor.rotatePhysical_akltRMPBasis,
        MPSTensor.akltTensorRMPPauli_neg_gaugeEquiv,
        MPSTensor.mpv_rotatePhysical_akltRMPBasis,
        MPSTensor.rotatePhysical_akltRMPBasisFixed}
    \leanok
    \uses{def:aklt_tensor_review, def:rotate_physical, thm:gauge_invariance,
        lem:mpv_smul, def:aklt_review_basis, lem:aklt_review_matrices, def:gauge_equiv}
    Let $U$ be the unitary whose columns are the vectors of
    (\ref{eq:aklt_review_basis}). The tensor (\ref{eq:aklt_review_tensor})
    expressed in this basis, that is, the tensor $B^{a}=\sum_{s}\overline{U_{sa}}\,A^{s}$
    obtained by rotating the physical index by $U^{\dagger}$, is
    \begin{align}
        B^{-} & = \tfrac{1}{\sqrt2}\sigma_x,
              & B^{+}                        & = -\tfrac{1}{\sqrt2}\sigma_y,
              & B^{0}                        & = \tfrac{1}{\sqrt2}\sigma_z.
        \notag
    \end{align}
    This is the form $\sigma_i/\sqrt2$ printed in
    \cite[Appendix~A, ``The AKLT state'']{Cirac2021Matrix} up to the virtual
    gauge $\sigma_y$ and the scalar $-1$: $B^{a}=\sigma_y\big(-\sigma_a/\sqrt2\big)\sigma_y^{-1}$
    for $a=-,+,0$ with $\sigma_-=\sigma_x$, $\sigma_+=\sigma_y$, $\sigma_0=\sigma_z$.
    Hence on $N$ sites the state coefficients of $B$ are $(-1)^N$ times those of
    the printed form, and the two describe the same state up to a global sign.
    In the basis with $\ket{+}=-\frac{i}{\sqrt2}(\ket{-1}+\ket{+1})$ the tensor
    becomes exactly $\sigma_i/\sqrt2$. The printed identity $A^{+}=\sigma_y/\sqrt2$
    in the basis (\ref{eq:aklt_review_basis}) holds only up to this gauge and sign,
    as recorded in \cite{gap:rmp_aklt_pauli_basis_sign}.
\end{theorem}
\begin{proof}
    \leanok
    Since $\langle +|\pm1\rangle=-i/\sqrt2$, one has
    $B^{+}=-\frac{i}{\sqrt2}(A^{-1}+A^{+1})=-\frac{i}{\sqrt2}Y=-\sigma_y/\sqrt2$,
    while $B^{-}=(A^{-1}-A^{+1})/\sqrt2=\sigma_x/\sqrt2$; conjugating the phase
    of $\ket{+}$ reverses the sign of $B^{+}$. Conjugation by $\sigma_y$ fixes
    $\sigma_y$ and reverses the signs of $\sigma_x$ and $\sigma_z$, which gives
    the gauge relation; a gauge leaves the state unchanged and the scalar $-1$
    multiplies it by $(-1)^N$.
\end{proof}

\begin{definition}[AKLT tensor]\label{def:aklt_tensor}
    \lean{MPSTensor.akltTensor}
    \leanok
    \uses{def:mps_tensor}
    Set $d = 3$ (spin-1) and $D = 2$. The AKLT tensor is the family
    $\{A^0, A^1, A^2\} \subset \MN{2}$ defined via the Pauli matrices by
    \begin{align}
        A^0 & = \frac{1}{\sqrt{3}} \sigma_z,
            & A^1                            & = \frac{\sqrt{2}}{\sqrt{3}} \ketbra{0}{1},
            & A^2                            & = -\frac{\sqrt{2}}{\sqrt{3}} \ketbra{1}{0}.
        \notag
    \end{align}
    Here $\sigma_z = \ketbra{0}{0} - \ketbra{1}{1}$, so equivalently
    $A^1 = \frac{1}{\sqrt{3}} \sigma_+$ and
    $A^2 = -\frac{1}{\sqrt{3}} \sigma_-$ for
    $\sigma_+ = \sqrt{2} \ketbra{0}{1}$ and
    $\sigma_- = \sqrt{2} \ketbra{1}{0}$. Equivalently, the matrices are
    the Clebsch--Gordan coefficients coupling spin-$1$ and
    spin-$\tfrac12$ to spin-$\tfrac12$. This is the representative of the
    tensor of Definition~\ref{def:aklt_tensor_review} used below
    (Lemma~\ref{lem:aklt_review_gauge_bridge}).
\end{definition}

\begin{lemma}[The singlet-bond tensor is a rescaled gauge of the representative]%
    \label{lem:aklt_review_gauge_bridge}
    \lean{MPSTensor.akltTensorRMP_gaugeEquiv}
    \leanok
    \uses{def:aklt_tensor_review, def:aklt_tensor, def:gauge_equiv, lem:aklt_review_matrices}
    With the physical labels $0,1,2$ of Definition~\ref{def:aklt_tensor} read as
    $S_z=0,+1,-1$, the tensor $A$ of Definition~\ref{def:aklt_tensor_review} and
    the tensor $T$ of Definition~\ref{def:aklt_tensor} satisfy
    $A^{s}=\sqrt{3/2}\,\sigma_z T^{s}\sigma_z^{-1}$ for every $s$.
\end{lemma}
\begin{proof}
    \leanok
    Conjugation by $\sigma_z$ fixes $\sigma_z$ and reverses the sign of
    $\ketbra{0}{1}$ and $\ketbra{1}{0}$; compare with
    Lemma~\ref{lem:aklt_review_matrices} entrywise.
\end{proof}

\begin{theorem}[AKLT tensor is not injective]\label{thm:aklt_not_injective}
    \lean{MPSTensor.aklt_not_isInjective}
    \leanok
    \uses{def:aklt_tensor, def:injective}
    The AKLT tensor is not injective at a single site:
    $\spn_{\C}\{A^0, A^1, A^2\}$ is the three-dimensional space of
    traceless matrices $\mathfrak{sl}(2,\C)$, which is a proper
    subspace of~$\MN{2}$.
\end{theorem}

\begin{theorem}[AKLT tensor is normal]\label{thm:aklt_normal}
    \lean{MPSTensor.aklt_isNormal}
    \leanok
    \uses{def:aklt_tensor, def:transfer_map}
    The transfer map of the AKLT tensor has a unique eigenvalue of
    maximal modulus (a simple spectral-radius eigenvalue); in
    particular the AKLT tensor is normal. Concretely, the
    length-$2$ blocked tensor is injective.
\end{theorem}

\begin{theorem}[$\Z_2 \times \Z_2$ on-site symmetry of the AKLT tensor]\label{thm:aklt_z2z2_symmetric}
    \lean{MPSTensor.aklt_isOnSiteSymmetric_Z2Z2}
    \leanok
    \uses{def:aklt_tensor, def:on_site_symmetric}
    The AKLT tensor is on-site symmetric under the $\Z_2 \times \Z_2$ subgroup
    of $\mathrm{SO}(3)$ generated by two commuting $\pi$-rotations about
    orthogonal axes, acting on the spin-$1$ physical index.
\end{theorem}

\begin{theorem}[Anticommuting virtual gauges of the AKLT symmetry]\label{thm:aklt_gauge_anticomm}
    \lean{MPSTensor.aklt_gauge_anticomm}
    \leanok
    The two virtual matrices that implement the symmetry of
    Theorem~\ref{thm:aklt_z2z2_symmetric}, namely $i\sigma_y$ and $\sigma_z$,
    anticommute.
\end{theorem}

\begin{definition}[AKLT factor system]\label{def:aklt_omega}
    \lean{MPSTensor.akltOmega}
    \leanok
    \uses{def:aklt_tensor}
    The AKLT factor system on $\Z_2 \times \Z_2$ is
    $\omega(g,h) = (-1)^{(g_1 + g_2)\,h_1}$ (the value $-1$ when
    $g_1 + g_2 = 1$ and $h_1 = 1$, and $1$ otherwise). It is the factor
    system of the explicit projective representation with virtual gauges
    $i\sigma_y$ and $\sigma_z$; the diagonal term $g_1 h_1$ relative to the
    cluster case reflects $(i\sigma_y)^2 = -\Id$.
\end{definition}

\begin{lemma}[AKLT commutator phase]\label{lem:aklt_comm_phase_neg_one}
    \lean{MPSTensor.aklt_commPhase_eq_neg_one}
    \leanok
    \uses{def:aklt_omega, def:comm_phase}
    The commutator phase of the AKLT factor system on the two generators is $-1$.
\end{lemma}

\begin{theorem}[The AKLT factor system is a $2$-cocycle]\label{thm:aklt_omega_cocycle}
    \lean{MPSTensor.akltOmega_isCocycle}
    \leanok
    \uses{def:is_cocycle, def:aklt_omega}
    The AKLT factor system is a genuine $2$-cocycle, being the factor system of an
    explicit projective representation, so it represents an element of
    $H^2(\Z_2 \times \Z_2, \mathrm{U}(1))$.
\end{theorem}

\begin{theorem}[The AKLT $\Z_2 \times \Z_2$ phase is non-trivial]\label{thm:aklt_nontrivial_spt}
    \lean{MPSTensor.aklt_isNontrivialSPT}
    \leanok
    \uses{def:nontrivial_class, thm:nontrivial_of_comm_phase, lem:aklt_comm_phase_neg_one}
    The AKLT factor system (Theorem~\ref{thm:aklt_omega_cocycle}), built from the
    anticommuting gauges $i\sigma_y$ and $\sigma_z$
    (Theorem~\ref{thm:aklt_gauge_anticomm}), is not a coboundary: its commutator
    phase on the two generators is $-1$. Its class in
    $H^2(\Z_2 \times \Z_2, \mathrm{U}(1)) \cong \Z_2$ is therefore non-trivial,
    so the AKLT chain realizes a non-trivial SPT phase rather than a trivial
    product state.
\end{theorem}

\begin{remark}\label{rem:aklt_z2z2_spt}
    The subgroup $\Z_2 \times \Z_2$ is the minimal subgroup of $\mathrm{SO}(3)$
    admitting a non-trivial $2$-cocycle.
\end{remark}

\begin{definition}[Blocked AKLT tensor]\label{def:aklt_blocked}
    \lean{MPSTensor.akltBlocked}
    \leanok
    \uses{def:aklt_tensor, def:blocked_tensor}
    The length-$2$ blocked AKLT tensor groups two physical sites into one, with
    blocked physical dimension $9 = 3^2$ and the same bond dimension $D = 2$.
\end{definition}

\begin{definition}[Blocked $\Z_2 \times \Z_2$ on-site action]%
    \label{def:aklt_blocked_z2z2_action}
    \lean{MPSTensor.akltBlockedZ2Z2Action}
    \leanok
    \uses{thm:aklt_z2z2_symmetric, def:blocked_tensor}
    The length-$2$ blocked on-site action of $\Z_2 \times \Z_2$ on the
    $9$-dimensional blocked physical space is the Kronecker square
    $g \longmapsto P_g \otimes P_g$ of the single-site spin-$1$
    $\pi$-rotation representation $P_g$ that implements the symmetry of
    Theorem~\ref{thm:aklt_z2z2_symmetric}. Identifying the blocked physical
    index with the pair of single-site indices, the entry of
    $P_g \otimes P_g$ at $(i_1 i_2, j_1 j_2)$ is the product
    $(P_g)_{i_1 j_1}(P_g)_{i_2 j_2}$ over the two grouped sites. This is the
    on-site action that makes the blocked twist agree with blocking the
    single-site twist.
\end{definition}

\begin{theorem}[The blocked AKLT tensor is injective]\label{thm:aklt_blocked_injective}
    \lean{MPSTensor.akltBlocked_isInjective}
    \leanok
    \uses{def:aklt_blocked, def:injective}
    The length-$2$ blocked AKLT tensor is injective: its nine matrices span
    $\MN{2}$.
\end{theorem}
\begin{proof}
    \leanok
    \uses{thm:aklt_normal}
    This is the blocked-tensor form of the normality of the single-site tensor
    (Theorem~\ref{thm:aklt_normal}): the single-site tensor is $2$-block
    injective, which is exactly injectivity of its length-$2$ blocking.
\end{proof}

\begin{theorem}[$\Z_2 \times \Z_2$ on-site symmetry of the blocked AKLT tensor]\label{thm:aklt_blocked_z2z2_symmetric}
    \lean{MPSTensor.aklt_blocked_isOnSiteSymmetric_Z2Z2}
    \leanok
    \uses{def:aklt_blocked, def:aklt_blocked_z2z2_action, def:on_site_symmetric, thm:aklt_z2z2_symmetric}
    The length-$2$ blocked AKLT tensor is on-site symmetric under
    $\Z_2 \times \Z_2$, with each group element acting by the Kronecker square
    $P_g \otimes P_g$ of the corresponding single-site spin-$1$ $\pi$-rotation
    $P_g$ (Definition~\ref{def:aklt_blocked_z2z2_action}).
\end{theorem}
\begin{proof}
    \leanok
    \uses{thm:aklt_z2z2_symmetric}
    Writing $A$ for the single-site tensor and $A^{[2]}$ for its length-$2$
    blocking, twisting the blocked tensor by $P_g \otimes P_g$ equals blocking
    the twisted single-site tensor,
    $(A^{[2]})^{P_g \otimes P_g} = (A^{P_g})^{[2]}$, where the superscript
    denotes the twist. Since the single-site tensor is symmetric, $A^{P_g}$ has
    the same matrix product vector family as $A$
    (Theorem~\ref{thm:aklt_z2z2_symmetric}), and blocking preserves the matrix
    product vector family, so $(A^{[2]})^{P_g \otimes P_g}$ has the same matrix
    product vector family as $A^{[2]}$.
\end{proof}

\begin{theorem}[The single-site $\Z_2 \times \Z_2$ action is unitary]\label{thm:aklt_unitary_z2z2}
    \lean{MPSTensor.aklt_isUnitary_Z2Z2}
    \leanok
    \uses{thm:aklt_z2z2_symmetric}
    Each single-site spin-$1$ $\pi$-rotation matrix $P_g$ of the
    $\Z_2 \times \Z_2$ action is unitary: $P_g P_g^\dagger = \Id$.
\end{theorem}

\begin{theorem}[The AKLT state has string order under its $\Z_2 \times \Z_2$ symmetry]\label{thm:aklt_string_order}
    \lean{MPSTensor.aklt_hasStringOrder}
    \leanok
    \uses{def:has_string_order, def:aklt_blocked, def:aklt_blocked_z2z2_action}
    For every element $g$ of $\Z_2 \times \Z_2$, the length-$2$ blocked AKLT
    tensor has string order in the sense of
    Definition~\ref{def:has_string_order} for the twist by $g$, with the
    maximally mixed boundary state $\Lambda = \tfrac{1}{2}\Id$ as its stationary
    boundary state.
\end{theorem}
\begin{proof}
    \leanok
    \uses{thm:has_string_order_of_symmetric_injective, thm:aklt_blocked_injective, thm:aklt_blocked_z2z2_symmetric, thm:aklt_unitary_z2z2, def:transfer_map}
    The blocked tensor is injective (Theorem~\ref{thm:aklt_blocked_injective}) and
    on-site symmetric under $\Z_2 \times \Z_2$
    (Theorem~\ref{thm:aklt_blocked_z2z2_symmetric}), and each group element acts
    by a real symmetric involutive matrix, hence unitarily. The maximally mixed
    state $\Lambda = \tfrac{1}{2}\Id$ is positive definite and has trace~$1$.
    The single-site AKLT matrices satisfy
    $\sum_i A^i(A^i)^\dagger = \Id$ and
    $\sum_i (A^i)^\dagger A^i = \Id$. These identities propagate to the
    length-$2$ blocked tensor $A^{[2]}$, whose matrices are the products
    $A^{i_1}A^{i_2}$. In particular, the blocked transfer map is unital:
    \begin{align}
        \E_{A^{[2]}}(\Id)
         & = \sum_{i_1,i_2}
        A^{i_1}A^{i_2}(A^{i_2})^\dagger(A^{i_1})^\dagger
        = \Id.
        \notag
    \end{align}
    Moreover, $\Lambda = \tfrac{1}{2}\Id$ is a fixed point of the adjoint
    transfer map:
    \begin{align}
        \E_{A^{[2]}}^\dagger(\Lambda)
         & = \sum_{i_1,i_2}
        (A^{i_2})^\dagger(A^{i_1})^\dagger
        \Lambda A^{i_1}A^{i_2}
        = \Lambda.
        \notag
    \end{align}
    The first identity sums the right-canonical relation
    $\sum_i A^i(A^i)^\dagger = \Id$ over the inner letter and then the outer;
    the second sums the left-canonical relation
    $\sum_i (A^i)^\dagger A^i = \Id$ over the outer letter and then the inner,
    using $\Lambda = \tfrac{1}{2}\Id$. Applying
    Theorem~\ref{thm:has_string_order_of_symmetric_injective} yields string
    order for every $g$.
\end{proof}

\begin{definition}[Cartesian AKLT tensor]\label{def:aklt_cartesian}
    \lean{MPSTensor.akltCartesian}
    \leanok
    \uses{def:mps_tensor}
    The Cartesian presentation of the AKLT tensor has $d = 3$, $D = 2$, and
    $A^x = \sigma_x$, $A^y = \sigma_y$, $A^z = \sigma_z$, the three Pauli
    matrices. This is the rotationally natural form used in
    \cite{Cirac2021Matrix} (around line~1159); it is the AKLT state in the
    Cartesian basis of the spin-$1$ site, the counterpart of the
    $\ket{m}$-basis form of Definition~\ref{def:aklt_tensor}.
\end{definition}

\begin{theorem}[The spin-$\tfrac12$ cover is surjective onto $\mathrm{SO}(3)$]\label{thm:su2_so3_surjective}
    \lean{SpinCover.spinHalfCover_surjective_onto_SO3}
    \leanok
    Every rotation $R \in \mathrm{SO}(3)$ is realized by conjugating the Pauli
    vector by some $U \in \mathrm{SU}(2)$: there is $U \in \mathrm{SU}(2)$ with
    $(\tfrac12\tr(\sigma_i U\sigma_j U^{-1}))_{ij} = R$.
    Equivalently, the adjoint double cover
    $\mathrm{SU}(2) \to \mathrm{SO}(3)$ is surjective.
\end{theorem}
\begin{proof}
    \leanok
    Factor $R = R_z(\alpha)R_x(\beta)R_z(\gamma)$ by its $ZXZ$ Euler
    decomposition and map each coordinate rotation to an explicit
    $\mathrm{SU}(2)$ generator:
    \begin{align}
        \diag(e^{-i\theta/2},e^{i\theta/2})
         & \mapsto R_z(\theta), \notag \\
        \begin{pmatrix}
            \cos(\beta/2)   & -i\sin(\beta/2) \\
            -i\sin(\beta/2) & \cos(\beta/2)
        \end{pmatrix}
         & \mapsto R_x(\beta).
        \notag
    \end{align}
\end{proof}

\begin{theorem}[$\mathrm{SO}(3)$ on-site symmetry of the AKLT tensor]\label{thm:aklt_so3_symmetric}
    \lean{MPSTensor.aklt_isOnSiteSymmetric_SO3}
    \leanok
    \uses{def:aklt_cartesian, def:on_site_symmetric, thm:su2_so3_surjective}
    The Cartesian AKLT tensor is on-site symmetric under the spin-$1$ (defining)
    representation of $\mathrm{SO}(3)$ on the physical index.
\end{theorem}

\begin{remark}\label{rem:aklt_so3_spt}
    The bond index carries the projective spin-$\tfrac12$ representation: the
    gauge for a rotation is either of its two $\mathrm{SU}(2)$ preimages, so the
    assignment is a representation only up to sign and factors through the
    double cover $\mathrm{SU}(2) \to \mathrm{SO}(3)$. Its class is the
    non-trivial element of
    $H^2(\mathrm{SO}(3),\mathrm{U}(1)) \cong \Z_2$, witnessing the SPT phase.
    % Maintainer note: the cohomological classification above (the class in
    % H^2(SO(3),U(1))) needs continuous group cohomology of SO(3) and is not
    % formalized; the obstruction is recorded in
    % docs/paper-gaps/rmp_aklt_continuous_rotation_gap.tex.
\end{remark}

\begin{theorem}[Symmetries and normality of the singlet-bond AKLT tensor]%
    \label{thm:aklt_review_symmetries}
    \lean{MPSTensor.akltTensorRMP_isNormal,
        MPSTensor.akltTensorRMP_isOnSiteSymmetric_Z2Z2,
        MPSTensor.akltTensorRMPPauli_isOnSiteSymmetric_SO3}
    \leanok
    \uses{lem:aklt_review_gauge_bridge, thm:aklt_normal, thm:aklt_z2z2_symmetric,
        thm:aklt_so3_symmetric, def:normal, def:on_site_symmetric,
        lem:on_site_symmetric_smul_gauge}
    The tensor of Definition~\ref{def:aklt_tensor_review} is normal and on-site
    symmetric under the $\Z_2\times\Z_2$ action of
    Theorem~\ref{thm:aklt_z2z2_symmetric}. Its Pauli form
    $\sigma_i/\sqrt2$ of Theorem~\ref{thm:aklt_review_pauli_form} is on-site
    symmetric under the spin-$1$ representation of $\mathrm{SO}(3)$.
\end{theorem}
\begin{proof}
    \leanok
    Normality and on-site symmetry are unchanged by a nonzero rescaling and by
    a virtual gauge (Lemma~\ref{lem:aklt_review_gauge_bridge}); the Pauli form
    is $1/\sqrt2$ times the Cartesian tensor of
    Definition~\ref{def:aklt_cartesian}.
\end{proof}

\begin{theorem}[The singlet-bond AKLT tensor lies in a non-trivial SPT phase]%
    \label{thm:aklt_review_spt}
    \lean{MPSTensor.akltTensorRMP_virtual_projRep}
    \leanok
    \uses{def:aklt_tensor_review, thm:aklt_z2z2_symmetric, thm:aklt_nontrivial_spt,
        def:aklt_omega, def:twisted_tensor}
    For every $g\in\Z_2\times\Z_2$, the twist of the tensor $A$ of
    Definition~\ref{def:aklt_tensor_review} by $g$ satisfies
    $\sum_j (P_g)_{ij}A^{j}\,X_g=X_g\,A^{i}$, where $X_g$ is the projective
    representation generated by $i\sigma_y$ and $\sigma_z$. Its factor system
    is the non-trivial class of $H^2(\Z_2\times\Z_2,\mathrm{U}(1))$
    (Theorem~\ref{thm:aklt_nontrivial_spt}).
\end{theorem}
\begin{proof}
    \leanok
    With the physical states ordered $S_z=0,+1,-1$, the two generators act by
    \begin{align}
        P_{(1,0)} & = \begin{pmatrix}-1&0&0\\0&0&1\\0&1&0\end{pmatrix},
                  & P_{(0,1)}                                           & = \begin{pmatrix}1&0&0\\0&-1&0\\0&0&-1\end{pmatrix},
        \notag
    \end{align}
    and the virtual matrices are
    \begin{align}
        X_{(1,0)} & = i\sigma_y = \begin{pmatrix}0&1\\-1&0\end{pmatrix},
                  & X_{(0,1)}                                            & = \sigma_z,
                  & X_{(1,1)}                                            & = X_{(1,0)}X_{(0,1)} = \begin{pmatrix}0&-1\\-1&0\end{pmatrix}.
        \notag
    \end{align}
    For $g=(1,0)$ the twist is $\widetilde A_g^{0}=-A^{0}$,
    $\widetilde A_g^{+1}=A^{-1}$, $\widetilde A_g^{-1}=A^{+1}$, and with the
    matrices of Lemma~\ref{lem:aklt_review_matrices},
    \begin{align}
        \widetilde A_g^{+1}X_g
         & = \ketbra{1}{0}\begin{pmatrix}0&1\\-1&0\end{pmatrix}
        = \ketbra{1}{1}
        = \begin{pmatrix}0&1\\-1&0\end{pmatrix}\big(-\ketbra{0}{1}\big)
        = X_g A^{+1},
        \notag                                                  \\
        \widetilde A_g^{0}X_g
         & = -\tfrac{1}{\sqrt2}\sigma_z\, i\sigma_y
        = \tfrac{1}{\sqrt2}\, i\sigma_y\,\sigma_z
        = X_g A^{0},
        \notag
    \end{align}
    since $\sigma_y$ and $\sigma_z$ anticommute; the letter $-1$ is checked in
    the same way. For $g=(0,1)$ the twist multiplies $A^{\pm1}$ by $-1$ and fixes
    $A^{0}$, which is conjugation by $\sigma_z$. The case $g=(1,1)$ is the
    composite of the two, and $g=1$ is trivial.
\end{proof}

\begin{theorem}[Two-site parent Hamiltonian of the singlet-bond AKLT tensor]%
    \label{thm:aklt_review_parent_hamiltonian}
    \lean{MPSTensor.akltTensorRMP_groundSpace_eq, MPSTensor.akltTensorRMP_mpvSubmodule_eq,
        MPSTensor.akltTensorRMP_chainGroundSpace_two_eq_mpvSubmodule,
        MPSTensor.akltTensorRMP_parentHamiltonian_two_unique_gs}
    \leanok
    \uses{lem:aklt_review_gauge_bridge, thm:aklt_chain_ground_space_two}
    The tensor of Definition~\ref{def:aklt_tensor_review} has the same local
    ground spaces, and a periodic matrix product state spanning the same line,
    as the representative of Definition~\ref{def:aklt_tensor}. Hence on a
    periodic chain of $N\geq3$ sites its two-site parent Hamiltonian has the
    AKLT state as its unique ground state.
\end{theorem}
\begin{proof}
    \leanok
    A virtual gauge and a nonzero rescaling leave the local ground spaces
    unchanged and multiply the periodic state by $(3/2)^{N/2}$; apply
    Theorem~\ref{thm:aklt_chain_ground_space_two}.
\end{proof}

%% ===================================================================
\subsection{Correlation length}\label{subsec:aklt_correlation_length}
%% ===================================================================

The AKLT state is the canonical example of a finite correlation length,
originally evaluated by Affleck, Kennedy, Lieb, and
Tasaki~\cite{AKLT1987Rigorous,AKLT1988ValenceBond} and Fannes, Nachtergaele, and
Werner~\cite{Fannes1992Finitely}. Its transfer map has leading eigenvalue $1$
and a single subleading eigenvalue $-\tfrac13$ of multiplicity three; the three
Pauli matrices are the subleading eigenvectors and the identity is the leading
one. The subleading modulus $\tfrac13$ gives the correlation length
$\xi = 1/\log 3$, and connected correlations decay as $(\tfrac13)^n$.

\begin{lemma}[Complex square of a real square root]
    \label{lem:complex_ofreal_sqrt_sq}
    \lean{Complex.ofReal_sqrt_sq}
    \leanok
    If $x\geq 0$, then the complex coercion of its real square root satisfies
    $((\sqrt{x}:\mathbb R):\mathbb C)^2=x$.
\end{lemma}

\begin{proof}
    \leanok
    Apply the real identity $(\sqrt{x})^2=x$ and commute the real embedding
    with powers.
\end{proof}

\begin{theorem}[Pauli eigenvectors of the AKLT transfer map]\label{thm:aklt_pauli_eigenvectors}
    \lean{MPSTensor.aklt_transferMap_sigmaZ, MPSTensor.aklt_transferMap_sigmaX,
        MPSTensor.aklt_transferMap_sigmaY}
    \leanok
    \uses{def:aklt_tensor, def:transfer_map}
    Each of the three Pauli matrices $\sigma_x$, $\sigma_y$, $\sigma_z$ is an
    eigenvector of the AKLT transfer map with eigenvalue $-\tfrac13$:
    \begin{align}
        \E_A(\sigma_k) & = -\tfrac13\,\sigma_k,
                       & k                      & \in \{x,y,z\}.
        \notag
    \end{align}
\end{theorem}

\begin{proof}
    \leanok
    \uses{lem:complex_ofreal_sqrt_sq}
    Each identity is an entrywise computation from the AKLT matrices, using
    $(\sqrt{3})^2 = 3$ and $(\sqrt{2})^2 = 2$ to clear the prefactors.
\end{proof}

\begin{theorem}[Subleading eigenvalue of the AKLT transfer map]\label{thm:aklt_subleading_eigenvalue}
    \lean{MPSTensor.aklt_transferMap_hasEigenvalue_one,
        MPSTensor.aklt_transferMap_hasEigenvalue_neg_third,
        MPSTensor.aklt_subleading_norm,
        MPSTensor.aklt_transferMap_one}
    \leanok
    \uses{thm:aklt_pauli_eigenvectors}
    The value $1$ is the leading eigenvalue of the AKLT transfer map, with the
    identity $\Id$ as eigenvector, and $-\tfrac13$ is a subleading eigenvalue,
    with $\sigma_z$ as eigenvector. The subleading modulus is therefore
    $\tfrac13$.
\end{theorem}

\begin{theorem}[AKLT correlation length]\label{thm:aklt_correlation_length}
    \lean{MPSTensor.aklt_correlationLength_eq, MPSTensor.aklt_correlationLength_pos}
    \leanok
    \uses{def:correlation_length, thm:aklt_subleading_eigenvalue}
    The AKLT correlation length, computed from the subleading eigenvalue
    $\lambda_2 = -\tfrac13$, equals
    \begin{align}
        \xi
         & = -\frac{1}{\log|\lambda_2|}
        = \frac{1}{\log 3}.
        \notag
    \end{align}
    It is finite and positive.
\end{theorem}

\begin{proof}
    \leanok
    \uses{lem:correlation_length_pos}
    The subleading modulus is $|{-\tfrac13}| = \tfrac13$, so rewriting the
    correlation length at $\lambda_2 = -\tfrac13$ gives
    $\xi = -1/\log(\tfrac13) = 1/\log 3$. Since $0 < \tfrac13 < 1$, positivity
    is the specialization of Lemma~\ref{lem:correlation_length_pos} to
    $\lambda_2 = -\tfrac13$.
\end{proof}

%% ===================================================================
